% ============================================================
%  DEOS Test Raporu -- Dokumantasyon
%  TEKNOFEST 2026 Robotaksi-Binek Ozerk Arac Yarismasi
% ============================================================
\documentclass[a4paper,12pt]{article}

\usepackage[utf8]{inputenc}
\usepackage[T1]{fontenc}
\usepackage{geometry}
\usepackage{amsmath}
\usepackage{amssymb}
\usepackage{booktabs}
\usepackage{longtable}
\usepackage{array}
\usepackage{xcolor}
\usepackage{listings}
\usepackage{hyperref}
\usepackage{enumitem}
\usepackage{fancyhdr}
\usepackage{multirow}
\usepackage{tcolorbox}
\usepackage{colortbl}
\usepackage{tabularx}

\geometry{
  left=2.5cm, right=2.5cm,
  top=2.5cm,  bottom=2.5cm
}

\hypersetup{
  colorlinks=true,
  linkcolor=blue!70!black,
  urlcolor=blue!60!black,
}

\definecolor{codebg}{rgb}{0.96,0.96,0.96}

\lstset{
  basicstyle=\ttfamily\footnotesize,
  keywordstyle=\color{blue!70!black}\bfseries,
  commentstyle=\color{green!50!black}\itshape,
  stringstyle=\color{red!60!black},
  numbers=left,
  numberstyle=\tiny\color{gray},
  numbersep=6pt,
  backgroundcolor=\color{codebg},
  frame=single,
  breaklines=true,
  tabsize=4,
  language=Python,
  showstringspaces=false,
  xleftmargin=1.8em
}

\tcbuselibrary{most}

\newtcolorbox{testbox}[2][]{
  colback=blue!4,
  colframe=blue!45!black,
  fonttitle=\bfseries\small,
  title=Test: \texttt{#2},
  #1
}

\newtcolorbox{specbox}[1][]{
  colback=green!5,
  colframe=green!45!black,
  fonttitle=\bfseries,
  title=Sartname Baglantisi,
  #1
}

\newtcolorbox{notebox}[1][]{
  colback=gray!6,
  colframe=gray!50!black,
  fonttitle=\bfseries,
  title=Not,
  #1
}

% Basit isaret komutlari (pifont gerektirmez)
\newcommand{\PASS}{\textcolor{green!60!black}{$\checkmark$}}
\newcommand{\FAIL}{\textcolor{red!70!black}{$\times$}}

\pagestyle{fancy}
\fancyhf{}
\rhead{DEOS Test Raporu}
\lhead{TEKNOFEST 2026}
\cfoot{\thepage}

\title{
  \vspace{-1cm}
  \textbf{\LARGE DEOS Test Raporu}\\[0.4em]
  \large Birim ve Entegrasyon Test Dokumantasyonu\\[0.2em]
  \normalsize TEKNOFEST 2026 Robotaksi-Binek Ozerk Arac Yarismasi\\[0.2em]
  \normalsize \texttt{test\_deos\_algorithms.py} --- \today
}
\author{BurakSahin00 / DEOS Takim}
\date{}

\begin{document}
\maketitle
\tableofcontents
\newpage

% ============================================================
\section{Test Altyapisi Genel Bakis}
% ============================================================

\subsection{Hedef}

Bu test dosyasi \texttt{deos\_algorithms} paketinin \textbf{davranissal
dogrulugunu} garantilemektedir. Her test:

\begin{itemize}
  \item Tek bir senaryo veya gereksinimi izole eder,
  \item Test adi kodlamayi bilmeyen biri tarafindan da anlasilabilir
        sekilde isimlendirilmistir,
  \item Sartnamede puanlanabilecek her kritik senaryo icin en az bir test
        bulunmaktadir.
\end{itemize}

\subsection{Test Yurutme Ortami}

\begin{lstlisting}[language=bash, caption=Test calistirma]
cd DEOS/deos_ws/src/deos_algorithms
python -m pytest -v --tb=short   # 120 test, ~0.32 s
\end{lstlisting}

\begin{center}
\begin{tabular}{ll}
  \toprule
  Cerceve      & \textbf{pytest} 7+\\
  Platform     & Windows 11 / Ubuntu 22.04 (ROS 2 Jazzy)\\
  Python       & 3.10+\\
  Toplam test  & \textbf{120}\\
  Calisma suresi & $\sim 0{,}32$ s\\
  Ortalama test basina & $\sim 2{,}7$ ms\\
  \bottomrule
\end{tabular}
\end{center}

\subsection{Test Organizasyonu}

Testler dosya icerisinde uc bolume ayrilmistir:

\begin{description}[leftmargin=2em]
  \item[Bolum 1] Temel fonksiyon testleri (93 test) ---
        her modulu bagimsiz olarak dogrular.
  \item[Bolum 2] Entegrasyon testleri ---
        birden fazla modul birlikte test edilir.
  \item[Bolum 3] Sartname bosluklari kapatma (27 test) ---
        TEKNOFEST 2026 sartnamesi incelenerek belirlenen eksik senaryolar.
\end{description}

\subsection{Paylasilan Yardimci Fonksiyonlar}

\begin{lstlisting}[language=Python, caption=Yardimci tespit yaraticilari]
def _obs_ped(dist_m, lateral_m=0.0, conf=0.92):
    """Yaya ObstacleDetection olusturur."""
    return ObstacleDetection(
        kind=ObstacleKind.PEDESTRIAN, confidence=conf,
        bbox_px=(280, 100, 360, 400),
        estimated_distance_m=dist_m,
        estimated_lateral_m=lateral_m,
    )

def _obs_cone(dist_m, lateral_m=0.0):
    """Koni ObstacleDetection olusturur."""
    ...

def _sign_det(class_name, dist_m=3.0, conf=0.95):
    """SignDetection olusturur."""
    return SignDetection(class_name=class_name,
                         confidence=conf,
                         bbox_px=(100, 50, 200, 150),
                         estimated_distance_m=dist_m)

def _confirm_sign(logic, det, now_base, frames=3):
    """CONFIRM_FRAMES kadar update cagirip son durumu dondurur."""
    st = None
    for i in range(frames):
        st = logic.update([det], now=now_base + i * 0.1)
    return st

def _slalom_cone_det(dist_m, lateral_offset):
    """Slalom testi icin koni olusturur."""
    cx = 320 + lateral_offset * 320   # 640px kamera
    ...
\end{lstlisting}

% ============================================================
\section{Bolum 1 --- Temel Fonksiyon Testleri}
% ============================================================

\subsection{Import ve Kurulum Testi}

\begin{testbox}{test\_imports}
Paketin tum alt modullerinin importlanabildigini dogrular.
Kurulum bozuksa (eksik dosya, sozdizim hatasi) bu test ilk patlar.

\smallskip
\textbf{Dogrulanan moduller:}
\texttt{safety\_logic}, \texttt{obstacle\_logic}, \texttt{slalom\_logic},
\texttt{parking\_logic}, \texttt{traffic\_sign\_logic},
\texttt{traffic\_light\_logic}, \texttt{geojson\_mission\_reader},
\texttt{waypoint\_manager}, \texttt{mission\_manager},
\texttt{perception\_fusion}, \texttt{sensors.types}
\end{testbox}

\subsection{SafetyLogic Testleri (9 test)}

\begin{testbox}{test\_safety\_logic\_yaya\_onay\_kareleri (parametrik)}
Ayni yaya BBox farkli kare sayilarinda sunulur.
\begin{itemize}
  \item 2 kare $\to$ \texttt{ThreatLevel.NONE} (onay yeterli degil)
  \item 3 kare $\to$ \texttt{ThreatLevel.HARD\_SLOW} (onay tamam)
\end{itemize}
\begin{specbox}
Sahte alarm onleme: $\geq 3$ ardisik kare onaylanmadan tehdit sayilmaz.
\end{specbox}
\end{testbox}

\begin{testbox}{test\_safety\_logic\_yanal\_esik}
Yanal mesafesi $> 1{,}2$ m (koridor disinda) olan nesne tehdit
sayilmamali: \texttt{ThreatLevel.NONE}.
\end{testbox}

\begin{testbox}{test\_safety\_logic\_korfez\_mesafe\_bandi (parametrik)}
Uc parametre: $d = 2{,}5$ m $\to$ EMERGENCY; $d = 6$ m $\to$ HARD\_SLOW;
$d = 12$ m $\to$ SOFT\_SLOW.
\end{testbox}

\begin{testbox}{test\_safety\_logic\_bbox\_mesafe\_tahmini}
Harici mesafe verilmezse bbox pin-hole formulu ile tahmin yapilmalidir.
Alt merkez pikseli $v$ ile mesafe formulu dogrulanir.
\end{testbox}

\begin{testbox}{test\_safety\_logic\_yanal\_offset\_pozitif\_sol}
Nesne goruntu merkezinin solundaysa $lateral\_m > 0$ olmalidir (isaret konvansiyonu).
\end{testbox}

\begin{testbox}{test\_safety\_logic\_yanal\_offset\_negatif\_sag}
Nesne goruntu merkezinin sagindaysa $lateral\_m < 0$ olmalidir.
\end{testbox}

\begin{testbox}{test\_safety\_logic\_harici\_mesafe\_oncelikli}
Harici mesafe ($estimated\_distance\_m$) verilmisse bbox tahmini
yerine harici deger kullanilmalidir.
\end{testbox}

\begin{testbox}{test\_safety\_logic\_caution\_carpani}
Yaya 6 m'de caution=1.5 ile etkin mesafe 4 m; koni ayni uzaklikta
etkin mesafe 6 m. Yaya daha yuksek tehdit uretmelidir.
\end{testbox}

\begin{testbox}{test\_safety\_logic\_tracker\_forget}
5 kare gorunmezse iz silinmeli ve tehdit uretilmemelidir.
\end{testbox}

\subsection{ObstacleLogic Testleri (7 test)}

\begin{testbox}{test\_obstacle\_logic\_empty}
Engel yokken: \texttt{speed\_cap=1.0}, \texttt{emergency\_stop=False}.
\end{testbox}

\begin{testbox}{test\_obstacle\_logic\_yaya\_yakinda\_bekleme}
Yaya $\leq 5$ m'de 5 kare sunulunca:
\texttt{waiting\_for\_dynamic\_obstacle=True}, \texttt{speed\_cap=0.0}.
\end{testbox}

\begin{testbox}{test\_obstacle\_logic\_yaya\_uzakta\_yavasla}
Yaya 8 m'de (SLOW bandinda): \texttt{behavior=DYNAMIC\_SLOW},
hiz tavan $\leq 0{,}4$.
\end{testbox}

\begin{testbox}{test\_obstacle\_logic\_koni\_serit\_degisimi}
Koni 2 m'de, 5 kare: \texttt{suggest\_lane\_change=True}.
\end{testbox}

\begin{testbox}{test\_obstacle\_logic\_ikiden\_fazla\_bariyer\_yol\_kapali}
2+ bariyer: \texttt{road\_blocked=True}.
\end{testbox}

\begin{testbox}{test\_obstacle\_logic\_dinamik\_bekleme\_sonra\_temizlendi}
Yaya kaybolunca DYNAMIC\_CLEAR\_FRAMES=3 kare beklenmeli, sonra
\texttt{waiting=False}.
\end{testbox}

\begin{testbox}{test\_obstacle\_logic\_statik\_serit\_degisimi\_bariyer}
Bariyer 2 m'de: \texttt{suggest\_lane\_change=True}.
\end{testbox}

\subsection{SlalomLogic Testleri (7 test)}

\begin{testbox}{test\_slalom\_no\_cones}
Koni yok: \texttt{faz=''bekleme''}, \texttt{aktif=False}.
\end{testbox}

\begin{testbox}{test\_slalom\_iki\_koni\_kapi\_baslangici}
Iki koni simetrik (sol/sag): \texttt{aktif=True},
\texttt{hedef\_taraf=''kapi''}.
\end{testbox}

\begin{testbox}{test\_slalom\_goruntu\_genisligi\_640px}
\textbf{D415 dogrulamasi:} \texttt{SlalomLogic().\_g\_w == 640},
\texttt{.\_g\_h == 480}. 1280 px olsaydi lateral offset yanlis hesaplanirdi.
\end{testbox}

\begin{testbox}{test\_slalom\_hiz\_katsayisi\_uzak / orta / yakin (parametrik)}
$> 3$ m $\to$ 1.0; $1{,}5$--$3$ m $\to$ 0.6; $< 1{,}5$ m $\to$ 0.3.
\end{testbox}

\begin{testbox}{test\_slalom\_tamamlanma\_15\_frame}
15 ardisik bos kare sonra \texttt{faz=''bitti''}.
\end{testbox}

\begin{testbox}{test\_slalom\_tek\_koni\_sagda\_sola\_steer}
Koni sagda ($offset > 0$) $\to$ $steering < 0$ (sola don).
\end{testbox}

\begin{testbox}{test\_slalom\_lateral\_offset\_640px\_dogru}
Merkez piksel (cx=320) $\to$ $lateral\_offset \approx 0{,}0$.
1280 px kamerada bu test $-0{,}5$ donup basarisiz olurdu.
\end{testbox}

\subsection{ParkingLogic Testleri (6 test)}

\begin{testbox}{test\_parking\_only\_allows\_parking\_sign\_spot}
\texttt{parking\_allowed=False} adayla manevra olmamali
(\texttt{no\_eligible\_spot=True}).
\end{testbox}

\begin{testbox}{test\_parking\_prefers\_allowed\_over\_closer\_forbidden}
Yasakli slot daha yakin olsa bile izinli slot secilmeli.
\end{testbox}

\begin{testbox}{test\_parking\_approach\_phase}
$d \leq 8$ m: \texttt{APPROACHING}; $d > 8$ m: \texttt{WAITING}.
\end{testbox}

\begin{testbox}{test\_parking\_align\_phase}
$d \leq 3$ m: \texttt{ALIGNING} fazina gecilmeli.
\end{testbox}

\begin{testbox}{test\_parking\_maneuver\_complete}
20 kare $d \leq 1{,}5$ m: \texttt{complete=True}.
\end{testbox}

\begin{testbox}{test\_parking\_search\_no\_detection}
Tespit yok: \texttt{no\_eligible\_spot=True}, arama steer uretilir.
\end{testbox}

\subsection{TrafficLightLogic Testleri (10 test)}

\begin{testbox}{test\_traffic\_light\_yellow\_ilkgorus\_yavaslama}
Ilk gorunen sari: \texttt{prepare\_to\_stop=True}, \texttt{cap=0.4}.
\end{testbox}

\begin{testbox}{test\_traffic\_light\_yellow\_duruyorken\_tavan\_sifir}
Sari + arac duruyor ($v \approx 0$): \texttt{cap=0.0}.
\end{testbox}

\begin{testbox}{test\_traffic\_light\_kirmizi\_sonra\_sari\_hazirlik}
Kirmizi onaylanmis $\to$ sari: \texttt{prepare\_to\_move=True},
\texttt{cap=0.25}.
\end{testbox}

\begin{testbox}{test\_traffic\_light\_kirmizi\_sonra\_sari\_duruyorken\_sifir}
Kirmizi $\to$ sari + duruyor: \texttt{prepare\_to\_move=True},
\texttt{cap=0.0}.
\end{testbox}

\begin{testbox}{test\_traffic\_light\_yesil\_sonra\_sari\_yavaslama}
Yesil $\to$ sari: \texttt{prepare\_to\_stop=True}, \texttt{cap=0.4}.
\end{testbox}

\begin{testbox}{test\_traffic\_light\_red\_emergency\_band\_must\_stop}
Kirmizi $d=2{,}5$ m ($\leq 3$ m acil bant):
\texttt{must\_stop=True}, \texttt{cap=0.0}.
\end{testbox}

\begin{testbox}{test\_traffic\_light\_red\_hard\_band\_slow}
Kirmizi $d=5$ m ($3$--$8$ m): \texttt{cap=0.5}.
\end{testbox}

\begin{testbox}{test\_traffic\_light\_red\_soft\_band\_slow}
Kirmizi $d=12$ m ($8$--$15$ m): \texttt{cap=0.8}.
\end{testbox}

\begin{testbox}{test\_traffic\_light\_red\_far\_no\_constraint}
Kirmizi $d=25$ m ($> 15$ m): \texttt{cap=1.0} (kisit yok).
\end{testbox}

\begin{testbox}{test\_traffic\_light\_red\_unknown\_soft\_when\_configured}
Mesafe bilinmiyor + \texttt{red\_unknown\_must\_stop=False}:
\texttt{cap=0.8} (soft yavas).
\end{testbox}

\subsection{TrafficSignLogic Testleri (7 temel test)}

\begin{testbox}{test\_traffic\_sign\_dur\_bekle}
\texttt{''dur tabelasi''} 5 sn boyunca \texttt{must\_stop\_soon=True} olmalidir.
\end{testbox}

\begin{testbox}{test\_traffic\_sign\_dur\_cooldown}
Dur tamamlandiktan sonra \texttt{STOP\_REARM\_SECONDS} icerisinde
ayni tabela tekrar tetiklenmemeli.
\end{testbox}

\begin{testbox}{test\_traffic\_sign\_girilmez}
\texttt{''girilmez''}: \texttt{must\_stop\_soon=True},
\texttt{entry\_blocked=True}.
\end{testbox}

\begin{testbox}{test\_traffic\_sign\_yaya\_gecidi}
\texttt{''yaya gecidi''}: hiz tavan $\leq 0{,}5$.
\end{testbox}

\begin{testbox}{test\_traffic\_sign\_tunel}
\texttt{''tunel''}: \texttt{approaching\_tunnel=True},
hiz tavan $\leq 0{,}7$.
\end{testbox}

\begin{testbox}{test\_traffic\_sign\_tur\_kisiti\_sol}
\texttt{''sol seridin sonu''} tur kisiti listesinde tanimli
(\texttt{active\_signs}'ta gorulmeli).
\end{testbox}

\begin{testbox}{test\_traffic\_sign\_kavsaktan\_sonra\_tur\_kisiti\_temizle}
\texttt{notify\_intersection\_passed()} sonrasi tur kisiti
\texttt{TurnPermissions} sifirlaniyor.
\end{testbox}

\subsection{WaypointManager Testleri (5 test)}

\begin{testbox}{test\_waypoint\_manager\_baslangic}
Plan yuklendi, ilk waypoint dogru sekilde donuyor.
\end{testbox}

\begin{testbox}{test\_waypoint\_manager\_varis}
Arac waypoint'e yetince \texttt{arrived=True}, wp\_index 1 artiyor.
\end{testbox}

\begin{testbox}{test\_waypoint\_manager\_bearing}
Hedef kuzey-dogu yonunde: bearing 0--90 arasi olmali.
\end{testbox}

\begin{testbox}{test\_waypoint\_manager\_cross\_track}
Yolun saginda olan arac: cross-track hatasi negatif olmali.
\end{testbox}

\begin{testbox}{test\_waypoint\_manager\_gorev\_tamamlandi}
Son noktayi gecince \texttt{mission\_complete=True}.
\end{testbox}

\subsection{MissionManager Testleri (4 test)}

\begin{testbox}{test\_mission\_manager\_pickup\_bekleme}
\texttt{PICKUP} noktasina varilinca: \texttt{speed\_cap=0.0},
\texttt{hold\_remaining\_s > 0}.
\end{testbox}

\begin{testbox}{test\_mission\_manager\_pickup\_tamamlandi}
15 s sonra \texttt{hold\_remaining\_s = 0}.
\end{testbox}

\begin{testbox}{test\_mission\_manager\_park\_modu}
\texttt{PARK\_ENTRY}: \texttt{park\_mode=True}, \texttt{park\_remaining\_s > 0}.
\end{testbox}

\begin{testbox}{test\_mission\_manager\_park\_zaman\_asimi}
180 s sonra \texttt{speed\_cap=0.0} (guvenli dur).
\end{testbox}

\subsection{GeoJsonMissionReader Testleri (4 test)}

\begin{testbox}{test\_geojson\_reader\_baslangic\_noktasi}
\texttt{''start''} adi veya gorevi: \texttt{TaskType.START}.
\end{testbox}

\begin{testbox}{test\_geojson\_reader\_park\_giris}
\texttt{''park\_giris''} adi: \texttt{TaskType.PARK\_ENTRY}.
\end{testbox}

\begin{testbox}{test\_geojson\_reader\_tumune\_crs}
\texttt{crs} alani varsa \texttt{plan.raw\_crs} doldurulmali.
\end{testbox}

\begin{testbox}{test\_geojson\_reader\_meta\_prefer\_tunnel}
\texttt{prefer\_tunnel=true} degerini dogru okumali.
\end{testbox}

\subsection{PerceptionFusion Testleri (2 test)}

\begin{testbox}{test\_perception\_fusion\_kisa\_senaryo}
Kamera+LiDAR girisi birlestirilince dogru ``kovalara'' ayrilmali:
kirmizi $\to$ isik; koni + lidar $\to$ engel; tabela $\to$ tabela.
\end{testbox}

\begin{testbox}{test\_parking\_detections\_from\_signs}
Park/park-yasak tabelasi $\to$ \texttt{ParkingDetection}:
\texttt{parking\_allowed=True/False}.
\end{testbox}

\subsection{DecisionArbiter Testleri (6 test)}

\begin{testbox}{test\_arbiter\_emergency\_stop\_baskın}
Herhangi bir aday \texttt{emergency\_stop=True} ise sonuc da
\texttt{True} olmali (OR kurali).
\end{testbox}

\begin{testbox}{test\_arbiter\_speed\_cap\_min}
Adaylar \texttt{[0.8, 0.4, 1.0]}: sonuc \texttt{cap=0.4} (MIN kurali).
\end{testbox}

\begin{testbox}{test\_arbiter\_steer\_oncelik (park $>$ dynamic)}
\texttt{park} (oncelik=0) ve \texttt{dynamic\_avoid} (oncelik=1):
park steer secilmeli.
\end{testbox}

\begin{testbox}{test\_arbiter\_steer\_dynamic\_avoid\_slalomdan\_ustun}
\texttt{dynamic\_avoid} (oncelik=1) ve \texttt{slalom} (oncelik=3):
dynamic\_avoid steer secilmeli.
\begin{lstlisting}[language=Python]
candidates = [
    Candidate(name="slalom",        steer_override=0.3),
    Candidate(name="dynamic_avoid", steer_override=-0.25),
]
assert dec.steer_override == pytest.approx(-0.25)
\end{lstlisting}
\end{testbox}

\begin{testbox}{test\_arbiter\_lane\_eksik\_avoidance\_iptal}
Serit bilgisi yok + zorunlu: avoidance steer iptal,
\texttt{LANE\_MISSING\_AVOIDANCE\_DISABLED} sebebi eklenir.
\end{testbox}

\begin{testbox}{test\_arbiter\_lane\_clamp\_steer\_kisaltir}
Dar serit: \texttt{steer=1.0} $\to$ clamp sonrasi $< 1{,}0$.
\end{testbox}

% ============================================================
\section{Bolum 3 --- Sartname Bosluklari Kapatma (27 Test)}
% ============================================================

\begin{specbox}
Bu testler TEKNOFEST 2026 Robotaksi-Binek sartnamesi incelenerek
belirlenmis bosluklar icin yazilmistir.
\end{specbox}

\subsection{ObstacleLogic Ek Senaryolari (4 test)}

\begin{testbox}{test\_obstacle\_logic\_dinamik\_kacinma\_yonu\_sol}
\textbf{Senaryo:} Yaya sagda ($lateral\_m=-0{,}5$), araç
dur-bekleme modundan cikiyor.

\textbf{Beklenti:} \texttt{dynamic\_avoidance\_direction=''left''}

\begin{lstlisting}[language=Python]
for i in range(20):
    state = logic.update(
        [_obs_ped(4.0, lateral_m=-0.5)], now=float(i)*0.05
    )
assert state.dynamic_avoidance_direction == "left"
\end{lstlisting}

\begin{specbox}
Tur-2/Tur-3: Dinamik engelden sakinma 50 puan. Gecis yonu
dogru olmalidir.
\end{specbox}
\end{testbox}

\begin{testbox}{test\_obstacle\_logic\_statik\_kacinma\_yonu\_sol\_konide\_saga}
\textbf{Senaryo:} Koni solda ($lateral\_m=+0{,}5$).

\textbf{Beklenti:} Sagdan gec: \texttt{avoidance\_direction=''right''}

Koordinat konvansiyonu: $lateral\_m > 0 \Rightarrow$ sol taraf
$\Rightarrow$ sagdan gecmek dogru hareket.
\end{testbox}

\begin{testbox}{test\_obstacle\_logic\_statik\_kacinma\_yonu\_sag\_konide\_sola}
\textbf{Senaryo:} Koni sagda ($lateral\_m=-0{,}5$).

\textbf{Beklenti:} Soldan gec: \texttt{avoidance\_direction=''left''}
\end{testbox}

\begin{testbox}{test\_obstacle\_logic\_statik\_commit\_stabilizasyon}
\textbf{Senaryo:} 3 kare sag koni (commit baslar), ardindan sol koni.

\textbf{Beklenti:} \texttt{STATIC\_AVOID\_COMMIT\_FRAMES} dolana kadar
yon degismemeli (zigzag onleme).

\begin{lstlisting}[language=Python]
for _ in range(3):
    state = logic.update([_obs_cone(2.0, lateral_m=0.5)])
committed_dir = state.avoidance_direction  # 'right'
state = logic.update([_obs_cone(1.8, lateral_m=-0.6)])
assert state.avoidance_direction == committed_dir  # degismemeli
\end{lstlisting}
\end{testbox}

\subsection{TrafficSignLogic Eksik Tabela Testleri (15 test)}

Sartname 22 tabela tanimlamaktadir. Onceki testlerde yalnizca 7 tabela
kapsaniyordu; 15 yeni test geri kalan tum tabelalar icin eklenmistir.

\begin{center}
\renewcommand{\arraystretch}{1.15}
\begin{longtable}{p{5.5cm}ll}
  \toprule
  Test Fonksiyonu                        & Beklenen Sonuc            & Puan\\
  \midrule
  \endhead
  \texttt{test\_...\_sola\_donulmez}     & \texttt{left=False}       & $-50$\\
  \texttt{test\_...\_saga\_mecburi}      & \texttt{forced=''right''} & $+50$\\
  \texttt{test\_...\_ileri\_mecburi}     & \texttt{forced=''straight''}& $+50$\\
  \texttt{test\_...\_ileri\_ve\_saga}    & \texttt{left=False}       & $+50$\\
  \texttt{test\_...\_ileri\_ve\_sola}    & \texttt{right=False}      & $+50$\\
  \texttt{test\_...\_ileriden\_saga}     & \texttt{forced=''right''} & $+50$\\
  \texttt{test\_...\_ileriden\_sola}     & \texttt{forced=''left''}  & $+50$\\
  \texttt{test\_...\_doner\_kavsak}      & \texttt{forced=''roundabout''}& $+50$\\
  \texttt{test\_...\_sagdan\_gidiniz}    & \texttt{forced=''pass\_right''}& $+50$\\
  \texttt{test\_...\_soldan\_gidiniz}    & \texttt{forced=''pass\_left''}& $+50$\\
  \texttt{test\_...\_park\_alani}        & \texttt{current\_area\_is\_parking}& $+80$\\
  \texttt{test\_...\_park\_yapilmaz}     & \texttt{current\_area\_no\_parking}& $-20$\\
  \texttt{test\_...\_isikli\_isaret}     & \texttt{traffic\_light\_expected}& ---\\
  \texttt{test\_...\_sol\_serit\_sonu}   & \texttt{active\_signs} icinde & $+50$\\
  \texttt{test\_...\_sag\_serit\_sonu}   & \texttt{active\_signs} icinde & $+50$\\
  \bottomrule
\end{longtable}
\end{center}

\begin{notebox}
\texttt{SignClass.KEEP\_RIGHT = ''sagdan gidiniz''} Turkce karakter icerir.
Test kodunda ham string yerine her zaman \texttt{SignClass.KEEP\_RIGHT}
sabitini kullanin.
\end{notebox}

\subsection{SlalomLogic Ek Senaryolari (2 test)}

\begin{testbox}{test\_slalom\_tek\_koni\_solda\_saga\_steer}
\textbf{Senaryo:} Koni solda (\texttt{lateral\_offset=-0.5},
$cx < 320$ piksel).

\textbf{Beklenti:} $steering > 0$ (saga don; koninin sagindan gec).

\begin{lstlisting}[language=Python]
st = sl.update([_slalom_cone_det(2.5, lateral_offset=-0.5)])
assert st.steering > 0.0
\end{lstlisting}

\begin{specbox}
640 px genis kamera: $offset = (cx - 320) / 320$.
$cx=160$ (sol) $\to$ $offset=-0{,}5$.
\end{specbox}
\end{testbox}

\begin{testbox}{test\_slalom\_faz\_zinciri\_bekleme\_aktif\_bitti}
Tam faz zinciri: bekleme $\to$ aktif $\to$ bitti (15 bos kare).

\begin{lstlisting}[language=Python]
assert sl.update([]).faz == "bekleme"
st = sl.update([_slalom_cone_det(2.0, lateral_offset=0.3)])
assert st.aktif is True
for _ in range(BITTI_FRAME_ESIK):
    state = sl.update([])
assert state.faz == "bitti"
\end{lstlisting}
\end{testbox}

\subsection{MissionManager Ek Senaryolari (2 test)}

\begin{testbox}{test\_mission\_manager\_dropoff\_hold\_tamamlandi}
\textbf{Senaryo:} DROPOFF noktasinda 14.9 s ve 16 s olcumu.

\begin{lstlisting}[language=Python]
_, dec_orta = mgr.update(pos, now_s=14.9)
assert dec_orta.hold_remaining_s > 0.0   # henuz bitmedi

_, dec_bitti = mgr.update(pos, now_s=16.0)
assert dec_bitti.hold_remaining_s == 0.0  # tamamlandi
\end{lstlisting}

\begin{specbox}
Sartname: Yolcu birakma bekleme suresi en az 15, en fazla 20 saniye.
Dogru bekleme: +70 puan.
\end{specbox}
\end{testbox}

\begin{testbox}{test\_mission\_manager\_pickup\_hold\_15\_20s\_penceresi}
15 s oncesinde hold devam etmeli; 16 s sonrasinda tamamlanmali.
Algoritma hicbir zaman 20 s'yi asmadigi da dogrulanir.
\end{testbox}

\subsection{GeoJsonMissionReader Ek Senaryolari (3 test)}

\begin{testbox}{test\_geojson\_reader\_gorev\_isimli\_noktalar}
\textbf{Senaryo:} GeoJSON'da \texttt{name=''gorev\_1''},
\texttt{''gorev\_2''}, \texttt{''gorev\_3''} adli noktalar.

\textbf{Beklenti:} Her biri \texttt{TaskType.CHECKPOINT} olmali.

\begin{specbox}
Gercek yaris GeoJSON'larinda gorev noktalari bu sekilde isimlendirilir.
Bu durum onceden kapsanmiyordu; checkpoint'lerin atlanma riski vardi.
\end{specbox}
\end{testbox}

\begin{testbox}{test\_geojson\_reader\_heading\_deg}
\texttt{heading\_deg=270} okunmali:
\texttt{plan.points[0].heading\_deg $\approx$ 270.0}.
\end{testbox}

\begin{testbox}{test\_geojson\_reader\_heading\_deg\_normalizasyon}
\texttt{heading\_deg=370} $\to$ $370 \bmod 360 = 10{,}0$ olmali.
\end{testbox}

\subsection{DecisionArbiter Ek Senaryosu (1 test)}

\begin{testbox}{test\_arbiter\_steer\_static\_avoid\_slalomdan\_ustun}
\textbf{Senaryo:} Slalom steer=0.3 ve static\_avoid steer=0.55 var.

\textbf{Beklenti:} \texttt{steer\_override = 0.55} (static\_avoid kazanmali).

\textbf{Aciklama:} \texttt{\_STEER\_PRIORITY}: static\_avoid=2, slalom=3.
Dusuk sayi = daha yuksek oncelik.

\begin{lstlisting}[language=Python]
candidates = [
    Candidate(name="slalom",       steer_override=0.3),
    Candidate(name="static_avoid", steer_override=0.55),
]
dec = arb.arbitrate(candidates=candidates, lane=valid_lane)
assert dec.steer_override == pytest.approx(0.55)
\end{lstlisting}
\end{testbox}

% ============================================================
\section{Test Kapsam Analizi}
% ============================================================

\subsection{Modul Bazli Test Sayisi}

\begin{center}
\begin{tabular}{lcc}
  \toprule
  Modul                               & Test Sayisi & Sartname Kapsamı\\
  \midrule
  \texttt{safety\_logic}              & 9           & Tum bant + tracker\\
  \texttt{obstacle\_logic}            & 11          & Dyn + statik + commit\\
  \texttt{traffic\_light\_logic}      & 11          & Tum bant + renk zinciri\\
  \texttt{traffic\_sign\_logic}       & 22          & 22/22 tabela sinifi\\
  \texttt{slalom\_logic}              & 9           & Tum hiz bant + faz\\
  \texttt{parking\_logic}             & 8           & Tum faz + slot secimi\\
  \texttt{waypoint\_manager}          & 8           & Bearing + XTE + varis\\
  \texttt{mission\_manager}           & 8           & Hold + park + timeout\\
  \texttt{geojson\_mission\_reader}   & 10          & Tum task + heading\\
  \texttt{decision\_arbiter}          & 7           & Oncelik + lane clamp\\
  \texttt{perception\_fusion}         & 3           & Fusion + park ayrimi\\
  Diger (import, route)               & 14          & Temel saglik\\
  \midrule
  \textbf{Toplam}                     & \textbf{120}& \textbf{Tam}\\
  \bottomrule
\end{tabular}
\end{center}

\subsection{Sartname -- Test Matrisi}

\begin{center}
\renewcommand{\arraystretch}{1.2}
\begin{longtable}{p{6.5cm}cr}
  \toprule
  Sartname Gereksinimi                   & Puan  & Kapsam\\
  \midrule
  \endhead
  Kirmizi isik ihlali                    & $-30$ & \PASS\\
  Kirmizi isikta dur (3 mesafe bandi)    & ---   & \PASS\\
  Yesil isik $\leq 5$ sn (+40)          & $+40$ & \PASS\\
  Yesil isik 5--30 sn (+20)             & $+20$ & \PASS\\
  Yesil isik $>30$ sn ($-20$)           & $-20$ & \PASS\\
  Sari isik yavaslama (4 senaryo)        & ---   & \PASS\\
  Tabela: 22 sinif                       & $\pm50$& \PASS\\
  Tabela: kavsak gecince temizle         & ---   & \PASS\\
  Tabela: DUR cooldown                   & ---   & \PASS\\
  Tabela: NO\_ENTRY entry\_blocked       & ---   & \PASS\\
  Tabela: TWO\_WAY hiz kisiti            & ---   & \PASS\\
  Yaya dinamik kacinma + yon             & $+50$ & \PASS\\
  Koni/bariyer statik kacinma + yon      & $+50$ & \PASS\\
  Zigzag onleme (commit stabilizasyon)   & ---   & \PASS\\
  Serit ihlali onleme (lane clamp)       & $-50$ & \PASS\\
  Serit yok: avoidance iptal             & ---   & \PASS\\
  Gorev noktasi radius dogrulama         & $+30$ & \PASS\\
  gorev\_1/2/3 isimleri                  & ---   & \PASS\\
  heading\_deg normalizasyon             & ---   & \PASS\\
  Pickup/Dropoff kafa acisi toleransi    & ---   & \PASS\\
  Pickup/Dropoff 15--20 s bekleme        & $+70$ & \PASS\\
  Park slotu secimi (izinli/yasakli)     & $+80$ & \PASS\\
  Park faz makinesi (5 faz)             & ---   & \PASS\\
  Park zaman asimi 180 s                 & ---   & \PASS\\
  Slalom koni sayimi + faz zinciri       & ---   & \PASS\\
  Slalom hiz bantlari (3 bant)           & ---   & \PASS\\
  Slalom steer yonu (sag + sol)          & ---   & \PASS\\
  640 px kamera kalibrasyon              & ---   & \PASS\\
  Tunel hiz indirimi                     & ---   & \PASS\\
  \bottomrule
\end{longtable}
\end{center}

\subsection{Basari Durumu}

\begin{center}
\begin{tabular}{lr}
  \toprule
  Olcut                 & Deger\\
  \midrule
  Toplam test sayisi    & 120\\
  Basarili              & 120 ($\%$100)\\
  Basarisiz             & 0\\
  Gecen sure            & $\sim 0{,}32$ s\\
  Ortalama test basina  & $\sim 2{,}7$ ms\\
  \bottomrule
\end{tabular}
\end{center}

% ============================================================
\section{Bilinen Kisitlamalar}
% ============================================================

\begin{enumerate}
  \item \textbf{End-to-end ROS testi yok:} Testler saf Python katmanini
        kapsar; ROS mesaj donusumu ve topic baglantilari bu dosyada
        dogrulanmaz.

  \item \textbf{GPS hassasiyeti:} \texttt{arrival\_radius\_m=3{,}0} m,
        UBLOX L86 GPS tipik dogrulugunu ($2$--$5$ m) karsilamak icin
        secilmistir. Gercek sahada bazen waypoint'e ucundan girilebilir.

  \item \textbf{Goruntu verisi:} Testler sentetik bbox degerleri kullanir;
        gercek YOLOv8 tespitlerinin gurultu orani veya aydinlatma
        degisimine gore farkli davranabilir.

  \item \textbf{Zaman belirleyicileri:} Tum zamanlayicilar
        \texttt{time.monotonic()} ile calisir; testlerde sahte
        \texttt{now} parametresi gecilir. Bu yaklasim CI ortaminda
        guvenilir tekrarlanabilirlik saglar.
\end{enumerate}

\end{document}
